\documentclass[12pt,a4paper]{article}

% ===== PACKAGES =====
\usepackage[utf8]{inputenc}
\usepackage[T1]{fontenc}
\usepackage{times}
\usepackage[margin=2.5cm]{geometry}
\usepackage{graphicx}
\usepackage{booktabs}
\usepackage{multirow}
\usepackage{array}
\usepackage{tabularx}
\usepackage{amsmath,amssymb}
\usepackage{enumitem}
\usepackage{hyperref}
\usepackage{xcolor}
\usepackage{float}
\usepackage{caption}
\usepackage{setspace}
\usepackage{titlesec}
\usepackage{fancyhdr}
\usepackage{natbib}

% ===== FORMATTING =====
\onehalfspacing
\setlength{\parindent}{0pt}
\setlength{\parskip}{6pt}

\hypersetup{
    colorlinks=true,
    linkcolor=blue!70!black,
    citecolor=blue!70!black,
    urlcolor=blue!70!black
}

\pagestyle{fancy}
\fancyhf{}
\rhead{\small HarnessEval --- Research Proposal}
\lhead{\small FPT University --- MSE-AI Program}
\cfoot{\thepage}

\titleformat{\section}{\Large\bfseries}{\thesection.}{0.5em}{}
\titleformat{\subsection}{\large\bfseries}{\thesubsection.}{0.5em}{}
\titleformat{\subsubsection}{\normalsize\bfseries}{\thesubsubsection.}{0.5em}{}

% ===== DOCUMENT =====
\begin{document}

% ===== TITLE PAGE =====
\begin{titlepage}
\centering
\vspace*{1cm}

{\large\textbf{FPT UNIVERSITY}}\\[0.3cm]
{\large Graduate School of Information Technology}\\[0.3cm]
{\normalsize Master of Software Engineering --- Artificial Intelligence}\\[2cm]

\rule{\textwidth}{1.5pt}\\[0.6cm]
{\LARGE\textbf{HarnessEval: Evaluating Coding Agent Scaffolds\\[0.3cm] Through Ablation-Driven Multi-Dimensional Analysis}}\\[0.6cm]
\rule{\textwidth}{1.5pt}\\[1.5cm]

{\large\textbf{Research Proposal}}\\[0.5cm]
{\normalsize Course: Methods of Learning and Scientific Research}\\[2cm]

\begin{tabular}{ll}
\textbf{Student:} & [Student Name] \\
\textbf{Student ID:} & [Student ID] \\
\textbf{Supervisor:} & [Supervisor Name] \\
\textbf{Date:} & April 2026 \\
\end{tabular}

\vfill
{\small Ho Chi Minh City, 2026}
\end{titlepage}

% ===== TABLE OF CONTENTS =====
\tableofcontents
\newpage

% ====================================================================
% CHAPTER 1: INTRODUCTION
% ====================================================================
\section{Introduction}

\subsection{Problem Statement}

Large Language Model (LLM)-based coding agents such as Claude Code, GitHub Copilot, Cursor, and SWE-Agent have become integral tools in modern software development. These agents autonomously read source code, fix bugs, generate features, and interact with file systems, terminals, and version control systems~\cite{A1,I1}.

However, the performance of a coding agent depends on two distinct factors: (1)~the \textbf{underlying LLM} (e.g., Claude, GPT, DeepSeek) and (2)~the \textbf{harness}---the infrastructure scaffold comprising the tool dispatch system, context management layer, and supporting components~\cite{A1}. Lou~et~al.~\cite{A2} demonstrated that the harness plays a critical role: a smaller model with a well-designed harness can outperform a larger model without one, eliminating 100\% of illegal actions in constrained environments. Bui~et~al.~\cite{A1} described the harness as an ``operating system for the agent,'' featuring a multi-layer architecture including dual-agent design, Adaptive Context Compaction (ACC) that reduces token usage by 54\%, and event-driven system reminders to combat instruction fade-out.

\textbf{The current paradox.} The research community has developed numerous benchmarks for evaluating agent \textit{output}---SWE-Bench~Pro with 1,865 tasks~\cite{C1}, SWE-Compass with 2,000 instances across 10 languages~\cite{C5}, and PRDBench with agent-driven evaluation~\cite{C4}. However, an analysis of 49 recent papers (2023--2026) reveals that \textbf{no study systematically evaluates the harness itself}~\cite{S5}. When reporting ``Agent~A achieves 70\% resolve rate,'' we do not know which harness component contributed to that result.

\textbf{Why this is a genuine research gap.} A legitimate question is: why have the groups that develop harnesses (Princeton---SWE-Agent, Anthropic---Claude Code) not conducted this evaluation themselves? Three reasons explain this gap:

\begin{enumerate}[leftmargin=2em]
    \item \textbf{Different roles.} They are \textit{developers} with incentives to build new harnesses, not \textit{evaluators} with incentives to compare systems---analogous to how automobile manufacturers build cars while independent organizations (Euro~NCAP, IIHS) conduct safety evaluations.
    \item \textbf{Rapid field growth.} From 2023 to 2026, the number of coding agents grew from a handful to dozens. No group has paused to meta-evaluate because the field is in a construction race.
    \item \textbf{Indirect evidence of need.} Lou~et~al.~\cite{A2} compared harnesses on game tasks; Robeyns~et~al.~\cite{B10} self-improved a harness without measuring individual component contributions---demonstrating that the need exists but has not been formalized.
\end{enumerate}

\textbf{Research gap.} There is no evaluation framework that enables: (a)~measuring the quality of individual harness components, (b)~decomposing the contributions of harness vs.\ LLM, and (c)~comparing different harnesses across unified quality dimensions.

\subsection{Research Purpose (Qualitative)}

This study aims to \textbf{construct the first integrated evaluation framework} for coding agent harness infrastructure, enabling the decomposition and quantification of individual harness component contributions to overall agent performance, independently of the underlying LLM capability.

\subsection{Research Objectives (Quantitative)}

\begin{description}[leftmargin=2em, style=nextline]
    \item[RO1.] Design a \textbf{3-dimensional harness evaluation taxonomy} (Tool Dispatch Efficiency, Context Utilization, Backend Portability) with 7~metrics total, validated by at least 5~domain experts (inter-rater agreement measured by Fleiss' $\kappa \geq 0.6$).

    \item[RO2.] Fork and refactor SWE-Agent into a \textbf{modular harness} with swappable components (tool set, context strategy, model backend), ensuring resolve rate \textbf{equivalent to baseline SWE-Agent} ($\pm 3\%$ absolute) on 50 verification tasks.

    \item[RO3.] Conduct a \textbf{systematic ablation study} on 150 tasks from SWE-Bench Verified with 27~conditions ($3 \times 3 \times 3$ full factorial: tool configurations $\times$ context strategies $\times$ LLM backends), including a pilot study (5~conditions $\times$ 20~tasks) prior to the full experiment. Report mean $\pm$ standard deviation from 3~runs for the 10 most critical conditions.

    \item[RO4.] Decompose variance contributions of harness vs.\ LLM using \textbf{three-way ANOVA} (Tool $\times$ Context $\times$ Backend), with statistical significance at $p < 0.05$ after Bonferroni correction. Include a \textbf{Generalized Linear Mixed Model} (GLMM) as a robustness check for binary outcomes.

    \item[RO5.] Release an \textbf{open-source evaluation toolkit} on GitHub with documentation, examples, and CI/CD pipeline.
\end{description}

\subsection{Research Hypotheses}

\begin{description}[leftmargin=2em]
    \item[H1.] The tool system has the \textbf{largest effect size} (Cohen's $d > 0.5$) among the three harness components on resolve rate.
    \item[H2.] Context management has a \textbf{medium effect size} (Cohen's $d = 0.3$--$0.5$).
    \item[H3.] Harness configuration explains at least \textbf{10\% of resolve rate variance} ($\eta^2 \geq 0.10$ in ANOVA), independent of LLM backend.
    \item[H4.] Resolve rate variance \textbf{across LLM backends decreases} when harness quality is higher (statistically significant interaction effect).
\end{description}

\subsection{Scope and Delimitations}

\textbf{In scope:}
\begin{itemize}[leftmargin=2em]
    \item Harness/scaffold infrastructure of LLM-based coding agents
    \item Three harness components: tool system, context manager, model backend
    \item 150 tasks from SWE-Bench Verified (Python)
    \item Three LLM backends: Claude Sonnet~4, GPT-4o, DeepSeek-V3
    \item Base harness: SWE-Agent (forked and refactored)
\end{itemize}

\textbf{Out of scope (future work):}
\begin{itemize}[leftmargin=2em]
    \item Safety/security evaluation $\rightarrow$ separate study (SecureCodeBench)
    \item Session continuity / long-horizon evaluation
    \item Languages other than Python
    \item Commercial harnesses (Claude Code, Cursor)
\end{itemize}

\subsection{Expected Thesis Structure}

\begin{enumerate}
    \item \textbf{Chapter~1:} Introduction (problem, literature review, objectives, scope)
    \item \textbf{Chapter~2:} Theoretical Background (agent architecture, ablation methodology, ANOVA)
    \item \textbf{Chapter~3:} Research Methodology (taxonomy design, modular harness, ablation protocol)
    \item \textbf{Chapter~4:} Results and Discussion
    \item \textbf{Chapter~5:} Conclusion and Future Work
\end{enumerate}

% ====================================================================
% CHAPTER 2: LITERATURE REVIEW AND THEORETICAL BACKGROUND
% ====================================================================
\section{Literature Review and Theoretical Background}

\subsection{Literature Review}

\subsubsection{Coding Agent Harness Architecture}

Wang~et~al.~\cite{S1} proposed the Profile-Memory-Planning-Action framework for LLM agents, cited over 3,000 times. Bui~et~al.~\cite{A1} instantiated this framework for coding agents with OpenDev---a terminal-native harness featuring dual-agent architecture (Planner-Executor), 5-stage Adaptive Context Compaction reducing token usage by 54\%, and event-driven system reminders. Lou~et~al.~\cite{A2} demonstrated that harnesses can be automatically synthesized, enabling smaller models to outperform larger ones. Cao~et~al.~\cite{A3} showed that coding agents process long contexts effectively by externalizing reasoning through file-system tools, outperforming SOTA by 17.3\%. Robeyns~et~al.~\cite{B10} proposed SICA---a self-improving coding agent that increased SWE-Bench from 17\% to 53\% over 15~iterations---but did not measure individual component contributions.

\textbf{Observation:} These works describe harness designs in detail but \textbf{do not benchmark} individual components.

\subsubsection{Coding Agent Benchmarks}

SWE-Bench Pro~\cite{C1} provides 1,865 contamination-resistant tasks from commercial codebases. SWE-EVO~\cite{C2} evaluates long-horizon software evolution (best model achieves only 21\%). SWE-Compass~\cite{C5} covers 10~programming languages and 8~task types. Prathifkumar~et~al.~\cite{C3} found evidence that LLMs may memorize patches rather than reason---Claude performs 3--6$\times$ better on file localization in SWE-Bench vs.\ novel benchmarks. AgentBoard~\cite{G3} introduced the progress rate metric with Pearson correlation $> 0.95$ with human judgment. A survey on agent evaluation~\cite{S5} found that 8 out of 10 popular benchmarks have serious validity issues---do-nothing agents pass 38\% of tasks.

\textbf{Observation:} All benchmarks evaluate \textbf{output quality}. None evaluates \textbf{harness quality}.

\subsubsection{Tool Learning and Context Management}

ToolLLM~\cite{F2} evaluates 16,000+ APIs with DFSDT reasoning. AnyTool~\cite{F3} proposes hierarchical self-reflective API retrieval. EasyTool~\cite{F5} reduces token cost by 70\%+ through documentation simplification. Mem0~\cite{E1} achieves 91\% latency reduction for persistent memory. A-MEM~\cite{E2} uses a Zettelkasten-inspired approach. Wang~et~al.~\cite{I2} discovered that agents repeat the same bugs across iterations due to lacking persistent memory.

\textbf{Observation:} Each paper evaluates \textbf{a single component in isolation}. No study evaluates \textbf{all components simultaneously} within a harness and compares their relative contributions.

\subsubsection{Research Gap Summary}

\begin{table}[H]
\centering
\caption{Research Gap Analysis: What Exists vs.\ What Is Missing}
\label{tab:gap}
\begin{tabularx}{\textwidth}{lXX}
\toprule
\textbf{Aspect} & \textbf{Existing Work} & \textbf{Gap (This Study)} \\
\midrule
Agent output evaluation & SWE-Bench, SWE-Compass, AgentBoard & --- \\
Individual tool accuracy & ToolLLM, AnyTool, Gorilla & Tool dispatch \textbf{within} coding harness \\
Individual memory evaluation & Mem0, A-MEM (chatbot only) & Memory \textbf{within} coding harness \\
Harness description & OpenDev, AutoHarness & \textbf{Measuring} harness quality \\
Cross-backend comparison & --- & \textbf{Entirely unexplored} \\
Harness vs.\ LLM decomposition & --- & \textbf{Entirely unexplored} \\
\bottomrule
\end{tabularx}
\end{table}

\subsection{Theoretical Background}

\subsubsection{Harness Architecture: Five-Layer Model}

Based on~\cite{S1} and~\cite{A1}, a coding agent harness comprises five layers:

\begin{enumerate}
    \item \textbf{Layer~1 --- Model Backend:} The underlying LLM (Claude, GPT, DeepSeek)
    \item \textbf{Layer~2 --- Tool System:} Available tools (Read, Write, Edit, Bash, Grep, Glob, etc.) and dispatch logic
    \item \textbf{Layer~3 --- Context \& Memory Management:} Context window management, compaction, persistent memory
    \item \textbf{Layer~4 --- Safety \& Permission:} Sandboxing, permission model, input validation (future work)
    \item \textbf{Layer~5 --- Orchestration:} Agent loop (ReAct, Plan-Execute), multi-turn management (future work)
\end{enumerate}

This study evaluates \textbf{Layers~1--3}. Layers~4--5 are designated as future work.

\subsubsection{Ablation Studies in Software Engineering}

Ablation study is a standard methodology in machine learning: systematically removing or modifying components to measure their individual contributions. Applied to harness evaluation:

\begin{itemize}[leftmargin=2em]
    \item Modify tool set (12 $\rightarrow$ 8 $\rightarrow$ 5 tools) $\rightarrow$ measure resolve rate change
    \item Change context strategy (full $\rightarrow$ sliding window $\rightarrow$ summary) $\rightarrow$ measure resolve rate change
    \item Switch LLM backend (Claude $\rightarrow$ GPT $\rightarrow$ DeepSeek) $\rightarrow$ measure portability
\end{itemize}

\textbf{Confounding variables.} When removing a tool, performance may decrease because (a)~the agent lacks the capability, or (b)~the harness dispatches poorly. To address this: we \textit{reduce} tools gradually rather than removing entirely, and record which tasks \textit{require} the removed tool to distinguish the two causes.

\subsubsection{Defining ``Independent of LLM'' via Three-Way ANOVA}

``Independent of LLM'' does \textbf{not} mean the harness and LLM do not interact. It means their contributions can be \textbf{statistically decomposed}.

\textbf{Method:} Three-way ANOVA (following reviewer feedback to use 3-way rather than 2-way):
\begin{itemize}[leftmargin=2em]
    \item Factor~A: Tool Configuration (3~levels)
    \item Factor~B: Context Strategy (3~levels)
    \item Factor~C: LLM Backend (3~levels)
    \item Dependent variable: Resolve rate
\end{itemize}

\begin{equation}
\text{Total Variance} = \underbrace{V_{\text{Tool}} + V_{\text{Context}}}_{\text{Harness}} + V_{\text{Backend}} + \underbrace{V_{\text{T} \times \text{B}} + V_{\text{C} \times \text{B}}}_{\text{Interactions}} + V_{\text{T} \times \text{C}} + V_{\text{T} \times \text{C} \times \text{B}} + V_{\text{Error}}
\end{equation}

If $V_{\text{Tool}}$ and $V_{\text{Context}}$ are statistically significant, the harness has an independent effect. The effect size $\eta^2$ indicates the percentage of variance explained by each factor. This formalization is the \textbf{primary theoretical contribution} of this study.

\textbf{Robustness check:} Since resolve rate per task is binary (0/1), ANOVA assumptions may be violated. A Generalized Linear Mixed Model (GLMM) with binary outcome and random effects for tasks will be used as a robustness check. Both ANOVA (for interpretability) and GLMM (for statistical correctness) will be reported.

% ====================================================================
% CHAPTER 3: RESEARCH METHODOLOGY
% ====================================================================
\section{Research Methodology}

\subsection{Overview}

This study employs an \textbf{experimental methodology} with four phases:
\begin{enumerate}
    \item Taxonomy design and expert validation
    \item SWE-Agent fork and modular refactoring
    \item Pilot study (5~conditions $\times$ 20~tasks)
    \item Full ablation study (27~conditions $\times$ 150~tasks)
\end{enumerate}

\subsection{Three-Dimensional Harness Evaluation Taxonomy (RO1)}

\begin{table}[H]
\centering
\caption{HarnessEval Taxonomy: 3~Dimensions, 7~Metrics}
\label{tab:taxonomy}
\small
\begin{tabularx}{\textwidth}{l l l X l}
\toprule
\textbf{Dim.} & \textbf{ID} & \textbf{Metric} & \textbf{Formal Definition} & \textbf{Auto?} \\
\midrule
\multirow{3}{*}{\shortstack[l]{D1. Tool\\Dispatch}}
    & M1.1 & Correct Selection & \% calls selecting an acceptable tool (from set of valid options, not single ground-truth) & No \\
    & M1.2 & Redundant Call Rate & \% calls whose output is unused within 3 subsequent turns & Yes \\
    & M1.3 & Utilization Breadth & Unique tool types used / total tools available & Yes \\
\midrule
\multirow{2}{*}{\shortstack[l]{D2. Context\\Utilization}}
    & M2.1 & Info Retention & BERTScore between full-context and compacted-context responses on same task & Yes \\
    & M2.2 & Effective Token Ratio & \% tokens classified as task-relevant by LLM classifier (validated on 200 segments with human labels, acceptance threshold: accuracy $>$ 90\%) & Semi \\
\midrule
\multirow{2}{*}{\shortstack[l]{D3. Backend\\Portability}}
    & M3.1 & Cross-Backend StdDev & $\sigma$(resolve rates across 3 backends) & Yes \\
    & M3.2 & Min/Max Ratio & $\min(\text{RR}) / \max(\text{RR})$ across backends & Yes \\
\bottomrule
\end{tabularx}
\end{table}

\textbf{Validation procedure:} The taxonomy will be presented to 5~domain experts (at least 2 from industry, at least 2 from academia). Each expert independently rates 20~sample evaluations. Inter-rater agreement will be measured using Fleiss' $\kappa$, with a target of $\kappa \geq 0.6$.

\subsection{Modular Harness via SWE-Agent Fork (RO2)}

SWE-Agent is an open-source (MIT license) coding agent with approximately 15,000~LOC. We fork and refactor three modules:

\begin{itemize}[leftmargin=2em]
    \item \texttt{tool\_config}: Swappable tool sets (12/8/5 tools) via configuration file
    \item \texttt{context\_strategy}: Swappable strategies (full/sliding-window/summary) via configuration file
    \item \texttt{model\_backend}: Swappable LLM backends (Claude/GPT/DeepSeek) via configuration file
\end{itemize}

\textbf{Acceptance criteria:} The modular harness must achieve resolve rate within $\pm 3\%$ absolute of baseline SWE-Agent on 50~verification tasks. The orchestration logic remains unchanged to ensure baseline quality.

\textbf{Estimated effort:} 2--3 weeks (vs.\ 1--3 months if built from scratch).

\subsection{Pilot Study (RO3, Phase~1)}

Prior to committing to the full experiment:

\begin{table}[H]
\centering
\caption{Pilot Study Design}
\label{tab:pilot}
\begin{tabular}{ll}
\toprule
\textbf{Parameter} & \textbf{Value} \\
\midrule
Conditions & 5 (Full, Reduced tools, No context mgmt, GPT backend, DeepSeek backend) \\
Tasks & 20 (randomly sampled from 150) \\
Runs per condition & 2 (for reproducibility estimate) \\
Total evaluations & $5 \times 20 \times 2 = 200$ \\
Estimated cost & \$50--80 \\
Duration & 1--2 weeks \\
\bottomrule
\end{tabular}
\end{table}

\textbf{Pilot objectives:}
\begin{enumerate}[leftmargin=2em]
    \item Validate that metrics discriminate between conditions
    \item Estimate accurate cost for the full experiment
    \item Debug the evaluation pipeline (Docker, API, logging)
    \item Compute preliminary effect sizes to confirm power analysis
\end{enumerate}

\subsection{Full Ablation Study (RO3, RO4, Phase~2)}

\begin{table}[H]
\centering
\caption{Full Factorial Experimental Design}
\label{tab:design}
\begin{tabular}{lll}
\toprule
\textbf{Factor} & \textbf{Levels} & \textbf{Details} \\
\midrule
Tool Configuration & 3 & Full (12 tools) / Medium (8) / Minimal (5) \\
Context Strategy & 3 & Full context / Sliding window (50K tokens) / Summary-based \\
LLM Backend & 3 & Claude Sonnet 4 / GPT-4o / DeepSeek-V3 \\
\midrule
\textbf{Total conditions} & \textbf{27} & $3 \times 3 \times 3$ full factorial \\
\bottomrule
\end{tabular}
\end{table}

\textbf{Dataset:} 150 tasks from SWE-Bench Verified, selected via stratified random sampling by difficulty (easy/medium/hard based on number of models that solved each task). Tasks with unstable Docker environments are excluded via pre-filtering.

\textbf{Runs:}
\begin{itemize}[leftmargin=2em]
    \item 10 most important conditions: 3~runs each (report mean $\pm$ std)
    \item 17 remaining conditions: 1~run each
    \item Total: $(10 \times 3 + 17 \times 1) \times 150 = 7{,}050$ evaluations
\end{itemize}

\textbf{Statistical analysis (RO4):}
\begin{itemize}[leftmargin=2em]
    \item \textbf{Three-way ANOVA:} Tool (3) $\times$ Context (3) $\times$ Backend (3)
    \item Report: $F$-statistic, $p$-value (Bonferroni-corrected), $\eta^2$ (eta-squared)
    \item Post-hoc: Tukey HSD for pairwise comparisons
    \item Effect size: Cohen's $d$ for each pairwise comparison
    \item \textbf{Robustness check:} GLMM with binary outcome (solved/not solved), random effects for tasks, fixed effects for Tool, Context, and Backend
    \item Power analysis: with 150 tasks, 27 conditions, $\alpha = 0.05$, power $\geq 0.80$
\end{itemize}

\subsection{Threats to Validity}

\subsubsection{Internal Validity}
\begin{itemize}[leftmargin=2em]
    \item \textbf{Confounding:} Reducing tools may change the problem, not just harness quality. \textit{Mitigation:} Reduce tools gradually (12$\to$8$\to$5); record which tasks require removed tools; exclude those from component analysis.
    \item \textbf{LLM non-determinism:} Stochastic outputs introduce variance. \textit{Mitigation:} Set temperature $= 0$; run 3~times for critical conditions; report confidence intervals.
    \item \textbf{Docker instability:} Some SWE-Bench environments are unstable. \textit{Mitigation:} Pre-filter tasks; retry twice on failure.
    \item \textbf{LLM versioning:} Models may be updated during the study. \textit{Mitigation:} Record exact model version/snapshot; run all experiments within a short time window.
\end{itemize}

\subsubsection{External Validity}
\begin{itemize}[leftmargin=2em]
    \item \textbf{Python only:} Tool systems for Python (pytest, pip) differ from Java (Maven, JUnit) or Rust (Cargo). Results require caution when generalizing.
    \item \textbf{Open-source only:} Commercial harnesses (Claude Code, Cursor) may exhibit different behavior.
    \item \textbf{Survivor bias:} Excluded tasks (Docker failures) may be harder tasks where harness matters more.
\end{itemize}

\subsubsection{Construct Validity}
\begin{itemize}[leftmargin=2em]
    \item \textbf{M1.1} depends on annotator judgment. \textit{Mitigation:} 2~annotators, Fleiss' $\kappa$.
    \item \textbf{M2.2} uses an LLM classifier. \textit{Mitigation:} Validate on 200~segments with human labels; acceptance threshold: accuracy $> 90\%$.
    \item \textbf{Statistical conclusion validity:} Multiple testing across 7~metrics requires Bonferroni or Benjamini-Hochberg FDR correction.
\end{itemize}

\subsection{Risk Analysis}

\begin{table}[H]
\centering
\caption{Risk Assessment and Mitigation}
\label{tab:risks}
\begin{tabularx}{\textwidth}{lllX}
\toprule
\textbf{Risk} & \textbf{Prob.} & \textbf{Impact} & \textbf{Mitigation} \\
\midrule
Docker failures on many tasks & Medium & Reduced sample & Pre-filter; backup tasks \\
API costs exceed budget & High & Cannot complete & Pilot estimates first; use DeepSeek (cheapest) \\
Modular harness $>$3\% below SWE-Agent & Low & RO2 failure & Minimal refactoring; keep original code \\
Metrics do not discriminate & Medium & Taxonomy invalid & Pilot detects early; revise metrics \\
Experts unavailable for validation & Medium & RO1 delayed & Online survey; increase expert pool \\
DeepSeek API instability & Low & Lose 1 backend & Backup: Qwen-2.5-72B or Llama-3.1-70B \\
\bottomrule
\end{tabularx}
\end{table}

% ====================================================================
% CHAPTER 4: EXPECTED RESULTS
% ====================================================================
\section{Expected Results}

\subsection{Expected Outcomes per Objective}

\textbf{RO1 --- Taxonomy:} A validated 3-dimensional, 7-metric taxonomy with Fleiss' $\kappa \geq 0.6$.

\textbf{RO2 --- Modular Harness:} A fork of SWE-Agent with modular components, resolve rate within $\pm 3\%$ of baseline.

\textbf{RO3 --- Ablation Study:} Empirical results answering:
\begin{itemize}[leftmargin=2em]
    \item Which component contributes most to resolve rate? (\textbf{H1}: Tool system has largest effect)
    \item How does context management affect performance? (\textbf{H2}: Medium effect size)
    \item What percentage of variance does harness explain? (\textbf{H3}: $\eta^2 \geq 0.10$)
    \item Does a better harness reduce cross-backend variance? (\textbf{H4}: Significant interaction)
\end{itemize}

Whether hypotheses are supported or rejected, both outcomes constitute scientific contributions.

\textbf{RO4 --- ANOVA Decomposition:} Variance decomposition table showing harness vs.\ LLM contributions.

\textbf{RO5 --- Toolkit:} Open-source Python package on GitHub with documentation and CI/CD.

\subsection{Scientific Significance}

\begin{itemize}[leftmargin=2em]
    \item \textbf{Theoretical:} Formalization of harness quality decomposition via ANOVA (first in this domain)
    \item \textbf{Empirical:} First quantification of individual harness component contributions
    \item \textbf{Methodological:} Reusable evaluation framework and toolkit
\end{itemize}

\subsection{Practical Significance}

\begin{itemize}[leftmargin=2em]
    \item \textbf{Developers:} Know which harness component to prioritize for improvement
    \item \textbf{Researchers:} Standardized framework for comparing new harnesses
    \item \textbf{Industry:} Evidence-based guidelines for harness design
\end{itemize}

\subsection{Publication Targets}

\begin{table}[H]
\centering
\caption{Target Publication Venues}
\begin{tabular}{lll}
\toprule
\textbf{Venue} & \textbf{Type} & \textbf{Fit} \\
\midrule
ICSE 2027 NIER & Conference (New Ideas) & Strong --- coding agent + evaluation \\
MSR 2027 & Conference & Good --- mining SE data \\
LLM Agents Workshop & Workshop (NeurIPS/ICML) & Good --- agent evaluation \\
EMNLP Demo Track & Conference & Good --- toolkit demonstration \\
\bottomrule
\end{tabular}
\end{table}

% ====================================================================
% CHAPTER 5: TIMELINE AND BUDGET
% ====================================================================
\section{Timeline and Budget}

\subsection{Timeline}

\begin{table}[H]
\centering
\caption{Research Timeline (24~Weeks)}
\label{tab:timeline}
\begin{tabularx}{\textwidth}{l l X l}
\toprule
\textbf{Phase} & \textbf{Weeks} & \textbf{Activities} & \textbf{Deliverable} \\
\midrule
GD1 & 1--6 & Taxonomy design + expert validation (5 experts, Fleiss' $\kappa$) & Taxonomy v1.0 \\
GD2 & 7--9 & Fork SWE-Agent + modular refactoring + verify $\pm 3\%$ & Modular harness v1.0 \\
Pilot & 10--11 & Pilot study (200 evaluations) & Pilot report \\
GD3 & 12--18 & Full ablation (7,050 evaluations) & Raw data, ANOVA \\
GD4 & 19--22 & Analysis + guidelines + thesis writing & Thesis draft \\
GD5 & 23--24 & Review, revision, toolkit release & Final thesis, GitHub \\
\bottomrule
\end{tabularx}
\end{table}

\subsection{Budget}

\begin{table}[H]
\centering
\caption{Budget Estimate}
\begin{tabular}{llr}
\toprule
\textbf{Item} & \textbf{Calculation} & \textbf{Cost (USD)} \\
\midrule
Pilot study & 200 evals $\times$ \$0.27 + retries & \$80 \\
Full experiment & 7,050 evals $\times$ \$0.27 + 30\% retries & \$2,500 \\
Contingency (20\%) & & \$500 \\
\midrule
\textbf{Total} & & \textbf{\$2,500--3,100} \\
\bottomrule
\end{tabular}
\end{table}

% ====================================================================
% REFERENCES
% ====================================================================
\newpage
\section*{References}
\addcontentsline{toc}{section}{References}

\begin{thebibliography}{99}

\bibitem{A1} N.~D.~Q. Bui~et~al., ``Building Effective AI Coding Agents for the Terminal: Scaffolding, Harness, Context Engineering, and Lessons Learned,'' \textit{arXiv:2603.05344}, 2026.

\bibitem{A2} X.~Lou~et~al., ``AutoHarness: Improving LLM Agents by Automatically Synthesizing a Code Harness,'' \textit{arXiv:2603.03329}, 2026.

\bibitem{A3} W.~Cao~et~al., ``Coding Agents are Effective Long-Context Processors,'' \textit{arXiv:2603.20432}, 2026.

\bibitem{B10} M.~Robeyns~et~al., ``A Self-Improving Coding Agent,'' \textit{arXiv:2504.15228}, 2025.

\bibitem{C1} X.~Deng~et~al., ``SWE-Bench Pro: Can AI Agents Solve Long-Horizon Software Engineering Tasks?'' \textit{arXiv:2509.16941}, 2025.

\bibitem{C2} M.~V.~T. Thai~et~al., ``SWE-EVO: Benchmarking Coding Agents in Long-Horizon Software Evolution Scenarios,'' \textit{arXiv:2512.18470}, 2025.

\bibitem{C3} T.~Prathifkumar~et~al., ``Does SWE-Bench-Verified Test Agent Ability or Model Memory?'' \textit{arXiv:2512.10218}, 2025.

\bibitem{C4} L.~Fu~et~al., ``Automatically Benchmarking LLM Code Agents through Agent-Driven Annotation and Evaluation,'' \textit{arXiv:2510.24358}, 2025.

\bibitem{C5} J.~Xu~et~al., ``SWE-Compass: Towards Unified Evaluation of Agentic Coding Abilities for Large Language Models,'' \textit{arXiv:2511.05459}, 2025.

\bibitem{E1} P.~Chhikara~et~al., ``Mem0: Building Production-Ready AI Agents with Scalable Long-Term Memory,'' \textit{arXiv:2504.19413}, 2025.

\bibitem{E2} W.~Xu~et~al., ``A-MEM: Agentic Memory for LLM Agents,'' \textit{arXiv:2502.12110}, 2025.

\bibitem{F2} Y.~Qin~et~al., ``ToolLLM: Facilitating Large Language Models to Master 16000+ Real-world APIs,'' \textit{arXiv:2307.16789}, ICLR 2024.

\bibitem{F3} Y.~Du~et~al., ``AnyTool: Self-Reflective, Hierarchical Agents for Large-Scale API Use,'' \textit{arXiv:2402.04253}, ACL 2024.

\bibitem{F5} S.~Yuan~et~al., ``EASYTOOL: Enhancing LLM-based Agents with Concise Tool Instruction,'' \textit{arXiv:2401.06201}, ACL Findings 2024.

\bibitem{G3} C.~Ma~et~al., ``AgentBoard: An Analytical Evaluation Board of Multi-Turn LLM Agents,'' \textit{arXiv:2401.13178}, NeurIPS 2024.

\bibitem{I1} ``AI Agentic Programming Survey,'' \textit{arXiv:2508.11126}, 2025.

\bibitem{I2} J.~Wang~et~al., ``Illuminating LLM Coding Agents: Visual Analytics for Deeper Understanding and Enhancement,'' \textit{arXiv:2508.12555}, IEEE TVCG 2025.

\bibitem{S1} L.~Wang~et~al., ``A Survey on Large Language Model based Autonomous Agents,'' \textit{arXiv:2308.11432}, Frontiers of Computer Science 2025.

\bibitem{S5} ``Survey on Evaluation of LLM-based Agents,'' \textit{arXiv:2503.16416}, 2025.

\end{thebibliography}

\end{document}
